\section{Quantitative positive localization on the unchanged source}
\label{sec:v63-positive-localization}
We apply the preceding moduli before pushforward. This preserves the
first-defect order, the complete incidence history, the image witness
$j+1$, and every exact label. The reference overcover defines a region;
all measure in that region still comes from the original Lorentz orbit.

For $0<\chi,h\le1$ let
$\widehat b^{\varepsilon,\mathrm{cau}(\chi,h),w}_{n,k,m,R}$ be the
first-clearance density of Section~\ref{sec:v48-physical-remainder},
restricted, before pushforward by the full roof, to
\begin{equation}\label{eq:v63-retained-source}
 c_i\ge\chi\quad(0\le i\le m),\qquad
 d_{m,\chi}<h.
\end{equation}
The selected witness is evaluated at $T_R^{j+1}x$, with its original
candidate disk and chart. The distance is
\eqref{eq:v63-buffered-distance}, using the cap $\chi/2$ for the
reference set only. For $w=1$ this is a positive restriction of the
same source, not a signed cancellation or a probability on
$\widehat{\mathcal K}$.

\begin{theorem}[Effective caustic localization at exact arithmetic labels]
\label{thm:v63-positive-localization}
Let $m\ge2$, $|w|\le M$, $0<2\varepsilon<s_0$ and $0<\chi,h\le1$.
Assume
\begin{equation}\label{eq:v63-source-threshold}
 4K_m\varepsilon\le(\chi h)^{E_m}.
\end{equation}
There are $C>0$, $A>1$, independent of these parameters, such that
\begin{align}\label{eq:v63-height-outside}
 &\operatorname*{ess\,sup}_t\sum_{n=1}^m\sum_{k\in\Z^2}
 \left|b^{\varepsilon,\mathrm{clr},w}_{n,k,m,R}(t)
       -\widehat b^{\varepsilon,\mathrm{cau}(\chi,h),w}_{n,k,m,R}(t)\right|
 \notag\\
 &\hspace{15mm}\le CMA^m\left((m+1)\chi+
                  K_m\varepsilon(\chi h)^{-E_m}\right).
\end{align}
For every bounded interval $J$ the retained uninserted density satisfies
\begin{equation}\label{eq:v63-mass-inside}
 \int_J m^2\widehat b^{\varepsilon,\mathrm{cau}(\chi,h),1}_{n,k,m,R}(t)\,dt
 \le C(1+|J|)^{144/145}
                  \left(K_m\chi^{-1}h^{1/E_m}\right)^{1/145}.
\end{equation}
The first assertion is essential-height control outside the retained
source. The second is mass control of that source, not its height.
\end{theorem}
\begin{proof}
Partition the actual first-clearance source into three positive pieces:
there is some $c_i<\chi$; all $c_i\ge\chi$ and $d<h$; all
$c_i\ge\chi$ and $d\ge h$. This is a partition on the original
trajectory, including incidences later than its first clearance defect.
On every regular trajectory,
\[
 \one_{\{\min_i c_i<\chi\}}
 \le\sum_{i=0}^m\chi c_i^{-1}\le(m+1)\chi\mathcal I_{m,R}.
\]
Theorem~\ref{thm:v61-inverse-incidence-density} pays the first piece
by $CMA^m(m+1)\chi$ in height. The middle piece is exactly
\eqref{eq:v63-retained-source}.

On the last piece Theorem~\ref{thm:v63-buffered-separation} and
\eqref{eq:v63-source-threshold} give
$|\mathcal J|\ge(\chi h)^{E_m}/(2K_m)$. Apply the physical area
formula and multiplicity bound of
Theorem~\ref{thm:v61-transverse-clearance} in the selected chart.
The marked source has the bounded rational-coordinate density inherited
from $\nu_R^*$, and the total length of the signed clearance bands,
summed over the geometrically weighted mark depths, is $O(\varepsilon)$.
The resulting height is at most
$CMA^mK_m\varepsilon(\chi h)^{-E_m}$.
Only for this positive upper comparison are exact-label restrictions
dropped. Restoring their disjoint partition at fixed $m$ gives the sum
in \eqref{eq:v63-height-outside}, with no target-count multiplier.
The source normalization contributes $c_*^{-1}$ once, absorbed in $C$.
For complex $w$, dominate its variation by $M$ times the same positive
source.

The normalized retained density is bounded above by the complete
normalized density $P=m^2p$. Its roof support is contained in the tube
of Theorem~\ref{thm:v63-effective-tube}. The layer-cake estimate of
Lemma~\ref{lem:v56-source-weight}, on a measurable $E\subset J$, is
\[
 \int_E P\le C(1+|J|)^{144/145}|E|^{1/145}.
\]
Insert the new tube bound to obtain \eqref{eq:v63-mass-inside}.
Coarea identities hold outside their usual Lebesgue-null roof sets;
no regular conditional version on an arithmetic zero class is used.
\end{proof}

\begin{corollary}[The raw-error formula with explicit cap and count budgets]
\label{cor:v63-effective-raw-error}
Use the original reconstruction width $\varepsilon=\varepsilon(B)$.
Under \eqref{eq:v63-source-threshold}, uniformly in $R,n,k$,
\begin{align}\label{eq:v63-effective-raw-error}
 \|P-G-m^2\widehat b^{\varepsilon,\mathrm{cau}(\chi,h),1}\|_\infty
 \le e_{B,m}+Cm^2A^m\left(\varepsilon+(m+1)\chi+
                     K_m\varepsilon(\chi h)^{-E_m}\right).
\end{align}
The reference $G$ is the unchanged signed canonical arithmetic kernel.
\end{corollary}
\begin{proof}
Substitute \eqref{eq:v63-height-outside} and the incidence-source
bound of Corollary~\ref{cor:v61-incidence-height-gain} into the positive
raw-error representation of Theorem~\ref{thm:v51-positive-error}.
The retained source is not discarded. No normalization of $G$ occurs
in this operation. In probability comparisons elsewhere, only its
positive part is normalized and probability TV is half the variation
mass; the sum in \eqref{eq:v63-height-outside} is variation mass itself.
\end{proof}

\paragraph{What the quantitative statement changes.}
For the buffered reference the fixed-count quantities in the earlier
separation and tube lemmas can now be replaced by
\[
 N=E_m,\qquad C_{m,\chi}=K_m\chi^{-E_m},\qquad
 D_{m,\chi}=K_m\chi^{-1},\qquad \beta=E_m^{-1}.
\]
These replacements concern the explicitly enlarged reference
\eqref{eq:v63-overcover} at cap $\chi/2$, not an asserted equality with
$\mathcal C_{m,\chi}$. The new estimate is uniform in the real cap;
it is stronger in this respect than a separate finiteness statement
for each $\chi$. Its exponent and prefactor follow from a fixed
quantifier-block budget and algebraic coefficient bounds. The finite
word geometry is not treated as a system of independent returns.

\paragraph{The ordered height problem.}
The budgets $E_m$ and $K_m$ are large. In particular they do not imply
\[
 \lim_{B\to\infty}\limsup_{m\to\infty}
 \sup_{R,n,k,\,\text{central }t}
        m^2b^{\varepsilon(B),\mathrm{inc}}_{n,k,m,R}(t)=0
\]
or its complete clearance analogue. In the required order, $B$ and
$\varepsilon(B)$ are fixed before the collision limit. Making the
finite-count error small by choosing $B$ after $m$ does not prove this
statement. Moreover \eqref{eq:v63-mass-inside} does not control the
height inside the tube. The two ordered heights in
Corollary~\ref{cor:v51-positive-criterion}, and hence its unrestricted
pointwise and exact-roof consequences, remain the next analytical
requirements on this same source.

\paragraph{A direct proof route for the physical source.}
The dependencies of the six geometric modules are as follows. Each row
separates the input from its conclusion; no height is inferred from mass alone.
\begin{center}
\begin{tabular}{@{}p{0.12\linewidth}p{0.39\linewidth}p{0.41\linewidth}@{}}
Module & Input on the actual orbit & Conclusion used downstream\\[3pt]
131 & Endpoint contact Hessian and coarea & Inverse-incidence density; incidence height\\
132 & Marked $(F,Z)$ and multiplicity & Transverse clearance height\\
133 & Full tangential jets on a roof fiber & Finite-type carriers, including folds\\
134 & Capped seam/rank graph & Finite caustic fibers and positive localization\\
135 & 134; fixed-format elimination & Buffered separation and tube budgets\\
136 & 131, 132, 135; positive raw error & Quantitative localization with its positive remainder
\end{tabular}
\end{center}
The determinant, physical graph and jet constructions remain subject
to independent specialist verification. Quantifier elimination supplies
moduli once those inputs have been established; source hashes and
finite fixtures do not verify the continuum geometry.
